\documentclass[11pt]{article}

% ===== arXiv-style preprint, monochrome =====
\usepackage[letterpaper,margin=1.1in]{geometry}
\usepackage[T1]{fontenc}
\usepackage[utf8]{inputenc}
\usepackage{amsmath,amssymb,amsthm}
\usepackage[round,authoryear]{natbib}
\usepackage[sc]{mathpazo}          % Palatino text + matching math, real small caps
\usepackage[scaled=0.88]{helvet}   % Helvetica for headings
\usepackage{microtype}
\usepackage{booktabs}
\usepackage{graphicx}
\usepackage{caption}
\usepackage{etoolbox}
\usepackage{fancyhdr}
\usepackage{xcolor}
\usepackage{listings}

% --- spacing / paragraphs ---
\setlength{\parindent}{1.2em}
\setlength{\parskip}{0.4em}
\linespread{1.08}
\setlength{\emergencystretch}{3em} % absorb long monospace identifiers

% --- hyperref: monochrome ---
\usepackage[hidelinks,breaklinks=true]{hyperref}
\usepackage{url}
\urlstyle{same}

% --- headings: sans, bold ---
\makeatletter
\renewcommand\section{\@startsection{section}{1}{\z@}%
  {-3.25ex \@plus -1ex \@minus -.2ex}{1.4ex \@plus.2ex}%
  {\normalfont\sffamily\Large\bfseries}}
\renewcommand\subsection{\@startsection{subsection}{2}{\z@}%
  {-2.6ex\@plus -1ex \@minus -.2ex}{1.0ex \@plus .2ex}%
  {\normalfont\sffamily\large\bfseries}}
\renewcommand\subsubsection{\@startsection{subsubsection}{3}{\z@}%
  {-2.2ex\@plus -1ex \@minus -.2ex}{0.8ex \@plus .2ex}%
  {\normalfont\sffamily\normalsize\bfseries}}
\makeatother
\setcounter{secnumdepth}{3}

% --- running header / footer ---
\pagestyle{fancy}
\fancyhf{}
\renewcommand{\headrulewidth}{0.4pt}
\renewcommand{\footrulewidth}{0pt}
\fancyhead[L]{\footnotesize\scshape AXL: Proof-Carrying Agent Spending Bounds}
\fancyhead[R]{\footnotesize\thepage}
\fancyfoot[C]{\footnotesize\itshape Preprint --- not peer reviewed}
\setlength{\headheight}{14pt}
\fancypagestyle{plain}{\fancyhf{}%
  \fancyfoot[C]{\footnotesize\itshape Preprint --- not peer reviewed}%
  \renewcommand{\headrulewidth}{0pt}}

% --- theorem environments ---
\theoremstyle{plain}
\newtheorem{theorem}{Theorem}
\newtheorem{lemma}{Lemma}
\newtheorem{proposition}{Proposition}
\theoremstyle{definition}
\newtheorem{definition}{Definition}
\theoremstyle{remark}
\newtheorem*{remark}{Remark}

% --- code listings ---
\lstdefinestyle{axl}{
  basicstyle=\ttfamily\small,
  breaklines=true,
  columns=fullflexible,
  frame=single,
  framesep=5pt,
  xleftmargin=4pt,
  showstringspaces=false,
  keepspaces=true,
}
\lstset{style=axl}

% --- title block ---
\makeatletter
\renewcommand{\maketitle}{%
  \begin{center}%
    {\LARGE\bfseries\@title\par}%
    \vspace{1.4em}%
    {\large\@author\par}%
    \vspace{0.35em}%
    {\normalsize Bluewave AI Research\par}%
    \vspace{0.15em}%
    {\normalsize\texttt{manuel@bluewaveai.online}\par}%
    \vspace{0.9em}%
    {\normalsize\@date\par}%
    \vspace{0.6em}%
    {\footnotesize\itshape
     Authored solely by Manuel Guilherme Galmanus. The work grew out of the
     author's agent-payment infrastructure in the Stellar ecosystem (Slippay)
     and was developed and shared in the context of the NEAR\,X ecosystem.\par}%
    \vspace{0.4em}%
  \end{center}%
  \vspace{0.6em}%
  \thispagestyle{plain}%
}
\makeatother

\renewenvironment{abstract}{%
  \vspace{0.4em}%
  \begin{center}\bfseries\normalsize Abstract\end{center}%
  \nopagebreak\par
  \begingroup
  \setlength{\leftskip}{2.2em}\setlength{\rightskip}{2.2em}%
  \small\noindent\ignorespaces
}{%
  \par\endgroup\vspace{1.0em}%
}

\hypersetup{
  pdftitle={AXL: A Standalone Proof-Carrying Compiler for Agent Spending Bounds},
  pdfauthor={Manuel Guilherme Galmanus}}

\title{AXL: A Standalone Proof-Carrying Compiler\\ for Agent Spending Bounds}
\author{Manuel Guilherme Galmanus}
\date{Preprint $\cdot$ June 2026}

\begin{document}
\maketitle

\begin{abstract}
Autonomous agents increasingly hold and move funds. The dominant safety
mechanism is a runtime check inside the agent's host or wallet --- but a
runtime check is exactly the thing a compromised host can skip, and a
counterparty has no way to confirm, before settling, that the check was ever
run. AXL is a small declarative language and a standalone Rust compiler
(\texttt{axlc}) that take a different stance: a spending policy is compiled
to a logical \emph{obligation}, discharged by an SMT solver via $k$-induction,
and emitted as a tamper-evident \emph{certificate} that any third party can
re-verify offline against the exact policy bytes. The compiler has zero crate
dependencies (standard library only; the solver is invoked as a subprocess and
its absence fails closed). For the deployed sliding-window policy we prove a
tight, minimal bound: total outflow in \emph{any} real-time window of length
$W$ is at most $2\times\text{window\_cap}$ --- the factor of two being a
provable consequence of epoch-boundary straddling, not a safety margin. We
describe the language (four orthogonal directives), the SMT lowering, the
certificate format and its SHA-256 binding, and a browser-side re-verifier
that regenerates the same obligations without a solver. A
continuous-integration \emph{gate} couples the proof to the deployment: on every
change touching the policy, the compiler, or the contract, CI re-discharges the
proof, verifies that the checked-in certificate reproduces, and asserts that the
constant the on-chain contract actually enforces equals the certified
multiplier --- so proof and deployed bound can no longer silently diverge (a
divergence turns CI red). We remain precise about residual scope: the contract
still enforces the bound by an in-source constant rather than reading the
certificate at runtime, and treats the policy hash as provenance; those are
stated as falsifiable conditions. The contribution is the verifiable artifact,
its construction, and the forcing function that keeps it honest under change.
\par\vspace{0.6em}
\noindent\textbf{Keywords:} proof-carrying code, SMT, $k$-induction, agent
autonomy, spending bounds, formal verification, smart wallets, AP2/x402.
\end{abstract}

\section{Introduction}

An autonomous agent that can pay needs a bound it cannot exceed. The standard
way to impose one is a runtime guard: the wallet or the orchestration host
checks, on each transaction, that the agent is within budget. This works until
the host is the adversary. A prompt-injected or compromised agent runtime can
route around its own guard; a counterparty receiving payment has no way to
inspect the payer's internal control and confirm, \emph{before} accepting
settlement, that any bound is in force. The check is real but unverifiable
from the outside, and unverifiable-from-the-outside is the wrong property for
money moving between mutually distrusting parties.

The classical answer to ``trust the producer's claim about code it ran'' is
proof-carrying code~\citep{necula1997}: the producer ships a machine-checkable
proof that its artifact satisfies a policy, and the consumer checks the proof
rather than trusting the producer. AXL applies that stance to agent spending.
A policy is written in a tiny declarative language; the compiler reduces the
quantitative part of the policy to a logical obligation; an SMT
solver~\citep{demoura2008} discharges it by $k$-induction~\citep{sheeran2000};
and the result is a small JSON \emph{certificate} cryptographically bound to
the exact policy text. Anyone --- a counterparty, an auditor, an insurer ---
can re-run the verification offline and either confirm the bound or detect that
the certificate has drifted from the policy.

\paragraph{What this paper claims, and what it does not.}
Honest positioning is load-bearing here. Agent \emph{payment} is not novel:
x402 carries machine payments over HTTP~402 and Google's AP2 standardizes
human-approved spending mandates. AXL is orthogonal to both --- it does not
move money; it proves a property of the policy that governs an agent that does.
The contribution is (i) a language whose quantitative directive lowers cleanly
to a decidable SMT obligation, (ii) a zero-dependency compiler that produces a
re-verifiable certificate with a tight, machine-checked bound, and (iii) a
faithful browser re-verifier, and (iv) a CI gate that re-discharges the proof
and asserts the contract's enforced multiplier equals the certified one on every
relevant change. We are equally clear about what remains:
\S\ref{sec:scope} states plainly that the contract still enforces the bound from
an in-source constant rather than reading the certificate at transaction time,
that the policy hash is provenance, and that the underlying SMT model is small
enough ($\approx$40 lines) that the proof alone is not a barrier to entry. We
treat the remaining items as falsifiable gaps to close, not as features to
obscure.

\paragraph{Contributions.}
\begin{itemize}
\item \textbf{AXL, a four-directive policy language} (\S\ref{sec:language})
  separating capability scope (\texttt{bind}), output schema
  (\texttt{constrain}), decidable predicates (\texttt{prove}), and the one
  SMT-discharged quantitative invariant (\texttt{invariant}). All four are
  enforced by code or solver, never by a language model.
\item \textbf{A standalone compiler} (\S\ref{sec:compiler}) with no crate
  dependencies: hand-rolled lexer, parser, JSON, and SHA-256, lowering the
  sliding-window invariant to SMT-LIB and discharging it via a subprocess
  solver, fail-closed when no solver is present.
\item \textbf{A tight, minimal, machine-checked bound} (\S\ref{sec:semantics}):
  for $\texttt{sliding\_window(ceiling}=2)$ the real-time-window outflow bound
  $K=2$ is proved inductive, attainable (tight), and minimal ($K=1$ is not
  inductive). The deployed wallet enforces exactly this constant.
\item \textbf{A proof-carrying certificate} (\S\ref{sec:certificate}) bound by
  SHA-256 to the policy bytes, with a defined load-bearing projection and
  drift detection, plus a browser re-verifier that regenerates the identical
  obligations without a solver (\S\ref{sec:verification}).
\item \textbf{An explicit scope-and-limits section} (\S\ref{sec:scope}) with
  falsifiable statements about what is and is not in the production path.
\end{itemize}

\section{The AXL Language}\label{sec:language}

An AXL program is an \texttt{agent} block with up to four orthogonal
directives. Each addresses a different layer of an agent's action, and each is
default-deny / fail-closed: an unrecognized construct is a compile error, never
a silent pass.

\begin{description}
\item[\texttt{bind -> [tools]}] declares the capability scope. A capability
  that is not bound has no schema in the action space, so a well-formed tool
  call for it cannot be emitted. An unknown capability raises
  \texttt{UnknownCapability} rather than being dropped silently.
\item[\texttt{constrain -> Schema}] fixes the output schema, compiled to
  constrained decoding (\texttt{json\_schema} with strict mode). A strict
  schema must declare \texttt{additionalProperties: false}, or the compiler
  refuses with \texttt{StrictSchemaNeedsAdditionalPropertiesFalse}.
\item[\texttt{prove -> <predicate>}] declares a decidable predicate over the
  output, of the form $\textit{lhs}\ \textit{op}\ \textit{rhs}$ with
  $\textit{op}\in\{<,\le,>,\ge,=,\ne,\texttt{in}\}$. Predicates are compiled to
  deterministic checks and evaluated at enforcement time; output violating any
  predicate is rejected.
\item[\texttt{invariant -> sliding\_window(ceiling = M) bound K}] is the only
  directive sent to the SMT solver. It declares a sliding-window spending
  policy with aggregate ceiling $M\cdot\text{window\_cap}$ and a
  \emph{claimed} real-time-window bound $K\cdot\text{window\_cap}$. Any other
  policy family (for example \texttt{token\_bucket}) is rejected with
  \texttt{UnsupportedPolicy} \emph{before} any solver runs.
\end{description}

\noindent The separation is deliberate: \texttt{bind} and \texttt{constrain}
shape what the agent \emph{can} express, \texttt{prove} checks decidable facts
about a specific output, and \texttt{invariant} is the one place where a
\emph{quantitative safety property over all possible action sequences} is
discharged by proof rather than by per-call checking.

\subsection{Worked examples}

The repository ships a matrix of examples exercising every branch of the
decision procedure. The deployed agent policy is the strict, sub-unity case
$M=2$:

\begin{lstlisting}
agent AgentWallet {
  bind      -> [read_balance, propose_payment]
  invariant -> sliding_window(ceiling = 2) bound 2
}
\end{lstlisting}

\noindent A non-invariant policy uses decidable predicates only:

\begin{lstlisting}
agent Settlement {
  bind      -> [read_balance, propose_settlement]
  constrain -> SettlementDecision
  prove     -> decision.amount <= account.balance
  prove     -> decision.recipient in account.allowlist
}
\end{lstlisting}

\noindent Table~\ref{tab:examples} summarizes the full matrix and the verdict
each yields. The two failing cases are first-class: an unbounded policy
(\texttt{ceiling = none}) and an unsupported family (\texttt{token\_bucket})
both \emph{refuse}, with distinct exit codes, so a missing or unsupportable
proof can never be mistaken for a passing one.

\begin{table}[t]
\centering
\resizebox{\linewidth}{!}{%
\small
\begin{tabular}{@{}llll@{}}
\toprule
Example & Policy & Verdict & Exit \\
\midrule
\texttt{agent\_wallet\_m1} & \texttt{sliding\_window(ceiling=1) bound 1} & ISSUED, $K{=}1$, tight & 0 \\
\texttt{agent\_wallet\_m2} & \texttt{sliding\_window(ceiling=2) bound 2} & ISSUED, $K{=}2$, tight & 0 \\
\texttt{agent\_wallet\_m3} & \texttt{sliding\_window(ceiling=3) bound 2} & ISSUED, $K{=}2$, tight + diagnostic & 0 \\
\texttt{agent\_wallet\_none} & \texttt{sliding\_window(ceiling=none) bound 2} & REFUSED (unbounded) & 2 \\
\texttt{agent\_wallet\_unsupported} & \texttt{token\_bucket(ceiling=2) bound 2} & REFUSED (unsupported) & 1 \\
\bottomrule
\end{tabular}%
}
\caption{The example matrix. The $M{=}3$ case is instructive: the weighted
throughput check binds tighter than the nominal ceiling, so the prover issues
$K{=}2$ and attaches a diagnostic that the declared ceiling is looser than the
proved bound.}
\label{tab:examples}
\end{table}

\section{The Compiler}\label{sec:compiler}

\texttt{axlc} is a standalone Rust crate (\texttt{axl-compiler} 0.1.0) with
\emph{no} \texttt{[dependencies]} section --- JSON, parsing, and SHA-256 are
implemented against the standard library only. The SMT solver is the single
external requirement, and it is a \emph{runtime} dependency invoked as a
subprocess, not a linked crate; its absence makes the prover fail closed rather
than silently pass. This keeps the trusted computing base of the verifier
small and auditable: a reader can read the entire compiler without chasing
transitive dependencies.

\subsection{Pipeline}

\begin{enumerate}
\item \textbf{Parse} ($\texttt{parse.rs}$): \texttt{parse\_agent} produces an
  \texttt{AgentSpec} with capabilities, schema, predicates, and an optional
  \texttt{InvariantDecl} (\texttt{family}, \texttt{ceiling}, \texttt{bound}),
  where the ceiling is a typed \texttt{Multiplier(u64)} or \texttt{None}.
\item \textbf{Compile} ($\texttt{compile.rs}$): \texttt{compile\_agent}
  resolves the toolset and schema into an \texttt{InferenceContract} and emits
  the request configuration (constrained decoding).
\item \textbf{Predicates} ($\texttt{predicate.rs}$):
  \texttt{compile\_predicate} parses each \texttt{prove} directive into a
  \texttt{Predicate}\,$(\textit{lhs},\textit{op},\textit{rhs})$ evaluated with
  JavaScript-faithful comparison semantics.
\item \textbf{SMT emission} ($\texttt{smt.rs}$): for an invariant,
  \texttt{emit} lowers each proof obligation to SMT-LIB text.
\item \textbf{Discharge} ($\texttt{prove.rs}$): \texttt{detect\_backend}
  probes for a solver, \texttt{check} runs it, and \texttt{discharge} returns a
  \texttt{Certificate} (\texttt{Issued}\,/\,\texttt{Refused}).
\item \textbf{Certify} ($\texttt{certify.rs}$): \texttt{build\_certificate}
  assembles the JSON artifact and binds it by SHA-256 to the policy bytes;
  \texttt{verify} re-derives and compares.
\end{enumerate}

\subsection{Solver dispatch, fail-closed}

\texttt{detect\_backend} probes, in order, a \texttt{z3} binary on
\texttt{PATH} (fed SMT-LIB on stdin) and a \texttt{python3} interpreter with
the \texttt{z3-solver} package. Verdict parsing reads the \emph{last}
\texttt{sat}/\texttt{unsat}/\texttt{unknown} token from solver output, so
solver warnings cannot corrupt the result. If no backend is found,
\texttt{discharge} returns \texttt{NoSolver} and the caller refuses: the
absence of a checker is never a pass.

\section{The Sliding-Window Model and the Safety Property}\label{sec:semantics}

The quantitative core of AXL is a single safety invariant over a two-epoch
abstraction of a sliding-window rate limiter, discharged by $k$-induction.

\subsection{State and transition}

Let $W$ be the window length (seconds) and let $\text{cap}$ abbreviate
$\text{window\_cap}$. The state is a pair $(p,c)$ giving outflow in the
previous and current epoch. A charge of amount $a$ arriving \texttt{elapsed}
seconds into the current epoch transitions the state as follows
(\texttt{smt.rs}), writing $\textit{rolled} \equiv (\texttt{elapsed} \ge W)$:
\begin{align*}
p_1 &= \textit{rolled} \,?\, (\texttt{elapsed} < 2W \,?\, c : 0) : p
   &&\text{(carry / drop / keep)}\\
c_1 &= \textit{rolled} \,?\, 0 : c\\
\textit{eie} &= \textit{rolled} \,?\, 0 : \texttt{elapsed}
   &&\text{(elapsed-in-epoch)}\\
\textit{weighted\_prev} &= \big\lfloor p_1 \cdot (W - \textit{eie}) / W \big\rfloor
   &&\text{(throughput shaping)}\\
\textit{accept} &\equiv (\textit{weighted\_prev} + c_1 + a \le \text{cap})
   \;\wedge\; (p_1 + c_1 + a \le M\cdot\text{cap})\\
c_2 &= c_1 + a.
\end{align*}
The acceptance test combines a \emph{weighted throughput} check (which smooths
spending across the boundary) with a hard \emph{aggregate ceiling}
$M\cdot\text{cap}$.

\subsection{The invariant and the obligations}

\begin{definition}[Window invariant]
For a claimed bound $K$,
\[
\textit{inv}_K(p,c) \;\equiv\; 0\le p\le \text{cap}\;\wedge\;0\le c\le \text{cap}\;\wedge\;p+c \le K\cdot\text{cap}.
\]
\end{definition}

\noindent The compiler emits up to four obligations and checks each verdict
against its expected value; any deviation refuses:
\begin{itemize}
\item \textbf{Base}: $\neg\,\textit{inv}_K(0,0)$ is \texttt{unsat} (the initial
  state satisfies the invariant).
\item \textbf{Inductive}: $\textit{inv}_K(p,c)\wedge \textit{accept}\wedge
  \neg\,\textit{inv}_K(p_1,c_2)$ is \texttt{unsat} (every accepted charge
  preserves the invariant).
\item \textbf{Attainable}: $\textit{inv}_K(p,c)\wedge \textit{accept}\wedge
  (p_1+c_2 = K\cdot\text{cap})$ is \texttt{sat} (the bound is reachable, hence
  tight rather than loose).
\item \textbf{Unbounded} (only for \texttt{ceiling = none}): a witness
  $n\cdot\text{per\_tx} > 1000\cdot\text{cap}$ is \texttt{sat}, forcing refusal.
\end{itemize}

\begin{theorem}[Deployed bound]\label{thm:bound}
For the deployed policy $\texttt{sliding\_window(ceiling}=2)$, the invariant
$\textit{inv}_2$ is inductive under \textit{accept}: every accepted charge from
a state satisfying $\textit{inv}_2$ lands in a state satisfying $\textit{inv}_2$.
Consequently the total outflow in any real-time window of length $W$ is at most
$2\cdot\text{window\_cap}$. The bound $K=2$ is \emph{tight} (attainable) and
\emph{minimal} ($\textit{inv}_1$ is not inductive).
\end{theorem}

\begin{proof}[Proof (machine-discharged)]
The Base and Inductive obligations for $K=2$ are emitted as SMT-LIB and
discharged \texttt{unsat} by the solver; \texttt{unsat} on the negated
inductive step is exactly the statement that no accepted transition can break
$\textit{inv}_2$. The Attainable obligation is \texttt{sat}, exhibiting a
charge sequence reaching $p+c = 2\cdot\text{cap}$, so the bound is not loose.
The corresponding inductive obligation for $K=1$ is \emph{not} \texttt{unsat},
so $1$ is not a sound bound; hence $2$ is minimal. The real-window consequence
follows because two epochs of length $W$ overlap any real $W$-length interval
at most pairwise, so accumulated outflow over such an interval is bounded by
$p+c$. \end{proof}

\begin{remark}[The factor of two is structural, not slack]
The $2\times$ is a provable consequence of \emph{epoch-boundary straddling}:
an adversary spends up to $\text{cap}$ immediately before a boundary and up to
$\text{cap}$ immediately after, and both charges fall inside one real
$W$-length interval. The weighted-throughput term shapes \emph{when} spending
lands but cannot defeat the straddle; $2\cdot\text{cap}$ is therefore the
honest worst case for a single-window abstraction, and AXL proves it rather
than assuming a tighter number that does not hold. Disanalogy with a fixed
calendar-window limiter: there the per-window cap is exactly $\text{cap}$ by
construction, but it permits a $2\times$ burst across the reset instant
\emph{without} reporting it; AXL's sliding model surfaces the same $2\times$
explicitly and certifies it.
\end{remark}

\section{The Proof-Carrying Certificate}\label{sec:certificate}

\texttt{axlc certify} emits a deterministic JSON certificate that binds the
verdict to the exact policy bytes. An ISSUED certificate for the deployed
policy has the shape:

\begin{lstlisting}
{
  "kind": "axl-proof-certificate",
  "axl_version": "0.1.0",
  "spec_sha256": "<sha256 of the policy bytes>",
  "agent": "AgentWallet",
  "invariant": { "family": "sliding_window", "ceiling": 2, "bound": 2 },
  "verdict": "ISSUED",
  "tight": true,
  "diagnostic": null,
  "onchain": {
    "ssl_hash": "<= spec_sha256>",
    "window_cap_multiplier": 2,
    "claim": "real-time window outflow <= 2 * window_cap",
    "matches_deployed_enforcement": true
  },
  "backend": "z3 binary (PATH)"
}
\end{lstlisting}

\noindent The binding is a hand-rolled FIPS\,180-4 SHA-256 of the policy bytes:
editing the policy invalidates the certificate. A \emph{load-bearing
projection} --- \texttt{kind}, \texttt{axl\_version}, \texttt{spec\_sha256},
\texttt{agent}, \texttt{family}, \texttt{ceiling}, \texttt{bound},
\texttt{verdict}, \texttt{tight}, \texttt{onchain.ssl\_hash}, and
\texttt{onchain.window\_cap\_multiplier} --- defines exactly which fields are
compared on re-verification; \texttt{backend} and the free-text \texttt{claim}
are metadata. A REFUSED certificate carries \texttt{verdict = "REFUSED"}, a
\texttt{refused\_reason}, and a null \texttt{onchain} block, so a refusal is a
first-class, signable outcome rather than an absent file.

\paragraph{What is, and is not, carried.}
The certificate carries the proved bound $K$, the policy SHA-256, and the
verdict --- not the SMT-LIB text and not the solver's internal proof object.
``Proof-carrying'' here means the bound claim is \emph{independently
re-derivable}: \texttt{verify-cert} re-emits the obligations from the
certificate's $(\textit{family},\textit{ceiling},\textit{bound})$, re-runs the
solver, and asserts the verdict and projection match. Drift between
certificate and policy is caught (exit 2); a certificate minted for one policy
does not verify against another, because the SHA-256 differs.

Disanalogy with classical proof-carrying code~\citep{necula1997}: there the
proof object itself travels and the consumer checks it directly with a fixed
proof checker; here the \emph{claim} travels and the consumer re-discharges the
obligation with its own solver. The trust shifts from ``check this proof'' to
``re-run this small, fully specified verification'' --- weaker if the consumer
lacks a solver, but the obligations are tiny and the policy binding is
cryptographic.

\section{Verification}\label{sec:verification}

\subsection{Offline browser re-verification}

\texttt{axlVerify.ts} re-verifies a certificate client-side with no solver
dependency. It (i) re-hashes the policy via WebCrypto and asserts it matches
\texttt{spec\_sha256} (cryptographic tamper-evidence); (ii) checks structural
coherence (\texttt{kind}, \texttt{verdict}, $\texttt{ssl\_hash} =
\texttt{spec\_sha256}$, $\texttt{window\_cap\_multiplier} = K$); and (iii)
\emph{regenerates} the exact SMT-LIB obligations \texttt{axlc} discharged, via
a TypeScript port of the emitter whose preamble, transition definitions, and
invariant expressions are byte-for-byte equivalent to \texttt{smt.rs}. The
browser thus confirms the binding and reproduces the theorems without running a
solver; discharging them automatically in-browser (e.g.\ via a WASM solver) is
identified as future work.

\subsection{Relationship to on-chain enforcement}

We state this precisely to avoid overclaiming. The Soroban smart-wallet
contract enforces the spending policy by \emph{hard-coded} arithmetic --- a
weighted throughput check and a hard ceiling
$\textit{prev\_spent} + \textit{cur\_spent} + \textit{amount} \le
2\cdot\text{window\_cap}$ --- and stores the policy hash (\texttt{ssl\_hash})
as immutable \emph{provenance only}. The contract does not read or verify any
certificate. The relationship between AXL and the contract is therefore
\emph{confirmatory}: AXL independently re-derives and machine-checks the same
$2\times$ constant the contract already encodes. AXL does not, today, feed the
bound to the contract, and no spend is cryptographically conditioned on the
proved policy.

\section{Evaluation}

The compiler is exercised by integration tests (\texttt{invariant\_prove.rs},
\texttt{certify.rs}) covering: parse-level rejection of malformed and
unsupported invariants; structural assertions on emitted SMT-LIB; end-to-end
solver discharge (skipped when no solver is installed); CLI exit codes; and
certificate round-trip, drift detection, and spec-swap rejection. The asserted
verdicts are exactly those of Table~\ref{tab:examples}: $M{=}2\to K{=}2$ tight,
$M{=}1\to K{=}1$ tight, $M{=}3\to K{=}2$ tight with a looser-ceiling
diagnostic, $\texttt{ceiling=none}\to$ REFUSED, and an overclaimed bound $\to$
REFUSED. The library is written to return \texttt{Result} on every parse and
JSON error and never to panic on malformed input. Concrete size: the SMT model
is on the order of $40$ lines of SMT-LIB; the compiler has $0$ external crate
dependencies; the certificate's load-bearing projection is $11$ fields.

\section{Scope and Limitations}\label{sec:scope}

AXL applies only to \emph{decidable} policy: arithmetic bounds, membership, and
bounded state. Irreducibly fuzzy judgments (``is this transaction
suspicious?'') are out of scope and still require a model in the loop. Within
that boundary, we describe what is now enforced under change and what residual
gap remains, each as a falsifiable statement.

\subsection{Closed: the proof is gated under change}\label{sec:gate}
A continuous-integration gate (\texttt{axl-compiler/scripts/axl-gate.sh},
wired in \texttt{.github/workflows/axl-gate.yml}) runs on every change touching
the policy, the compiler, or the contract. It (i) re-discharges the invariant
with z3 and asserts the checked-in certificate reproduces (\texttt{verify-cert},
exit~2 on drift), and (ii) extracts the multiplier the contract actually
enforces --- the integer $N$ in \texttt{window\_cap.saturating\_mul(}$N$\texttt{)}
--- and asserts $N$ equals the certified \texttt{window\_cap\_multiplier}. We
validated all three failure paths: with the repository clean the gate exits~0;
mutating the contract literal to $3$ makes it exit~2 (``contract enforces
$3\times$ \dots\ proof certifies $2\times$''); tampering the certificate's
multiplier makes \texttt{verify-cert} exit~2. \emph{Falsifiable:} a commit that
changes the deployed bound without re-proving it, or that re-proves a different
bound without updating enforcement, cannot reach \texttt{main} green. This is
the concrete sense in which the proof is now load-bearing under change rather
than an artifact that happens to agree with the contract today.

\begin{itemize}
\item \textbf{Residual: the contract reads a constant, not the certificate at
  runtime.} The enforced multiplier is an in-source literal; the gate binds it
  to the certified value at \emph{build} time, not at transaction time.
  \emph{Falsifiable:} a deployment that ships a contract binary whose literal
  was edited after the gate ran would not be caught on-chain --- closing this
  fully requires the contract (or the deploy step) to read the certified bound,
  which is a contract change we did not make autonomously.
\item \textbf{The policy hash is provenance, not a runtime gate.}
  \texttt{ssl\_hash} is pinned immutably but not interpreted by
  \texttt{\_\_check\_auth}. \emph{Falsifiable:} no individual spend is rejected
  on the basis of a hash mismatch with a certified policy.
\item \textbf{The proof is not a moat.} The SMT model is $\approx$40 lines; a
  competent formal-methods engineer reproduces it in engineer-weeks. The
  durable advantage is the \emph{gate plus adoption}: a certificate that is
  required-to-participate (a second prover, an insurer, an on-chain registry).
  The gate now exists; broad adoption does not yet.
\item \textbf{The solver is a runtime dependency.} \texttt{prove},
  \texttt{certify}, and \texttt{verify-cert} require z3 (binary or Python
  package). This is intentional and fail-closed --- absence refuses --- but a
  consumer without a solver can check only the binding and the obligation text,
  not the discharge.
\end{itemize}

\noindent We surface these because the value of a proof-carrying scheme is
exactly its honesty about what has and has not been checked. The certificate
machinery --- language, lowering, $k$-induction discharge, SHA-256 binding,
deterministic re-verification, drift detection --- is real, tested, and now
gated in CI. The remaining work is runtime, not build-time: making the contract
itself consume the certified bound.

\section{Related Work}

\textbf{Proof-carrying code}~\citep{necula1997} originates the
producer-ships-proof, consumer-checks-proof stance; AXL adapts it to a
re-discharge model bound by a content hash (\S\ref{sec:certificate}).
\textbf{SMT solving}~\citep{demoura2008} and \textbf{$k$-induction}
\citep{sheeran2000} are the engine and the proof technique for the
sliding-window invariant. \textbf{Smart-contract verification} (e.g.\
specification and model checking of on-chain logic) targets the contract
itself; AXL instead certifies the \emph{off-chain policy} an agent commits to,
and (today) only confirms the contract's hard-coded bound rather than
generating it. \textbf{Agent-payment standards} --- x402 for HTTP-native
machine payments and Google's AP2 for human-approved mandates --- provide the
rails and the authorization chain; AXL is orthogonal, proving a quantitative
property of the governing policy rather than transporting value or capturing
consent. \textbf{Capability security} informs the \texttt{bind} directive's
default-deny action-space scoping.

\section{Conclusion}

AXL reframes an agent's spending limit from a runtime check that a compromised
host can skip into a logical obligation that anyone can re-verify offline. The
compiler is small and dependency-free by construction, the sliding-window
bound is proved tight and minimal by $k$-induction, and the certificate is
bound to its policy by a content hash so that drift is detectable. A CI gate
now re-discharges that proof on every relevant change and refuses to let the
deployed multiplier diverge from the certified one, so the agreement between
proof and enforcement is checked rather than coincidental. The honest residual
is runtime, not build-time: the contract still enforces an in-source constant
rather than reading the certified bound at transaction time, and the policy hash
remains provenance. That is the falsifiable condition whose resolution would
make the scheme load-bearing end to end. The thesis is narrow and, we think,
correct: for money moving between agents that do not trust each other, a bound
the counterparty can re-check beats a bound it must take on faith.

\begin{thebibliography}{9}
\small
\bibitem[Necula(1997)]{necula1997}
Necula, G.~C. (1997). Proof-carrying code. In
\emph{Proc.\ 24th ACM SIGPLAN-SIGACT Symposium on Principles of Programming
Languages (POPL)}, 106--119.

\bibitem[de Moura \& Bj\o rner(2008)]{demoura2008}
de Moura, L., \& Bj\o rner, N. (2008). Z3: An efficient SMT solver. In
\emph{Tools and Algorithms for the Construction and Analysis of Systems
(TACAS)}, LNCS 4963, 337--340.

\bibitem[Sheeran et al.(2000)]{sheeran2000}
Sheeran, M., Singh, S., \& St\aa lmarck, G. (2000). Checking safety properties
using induction and a SAT-solver. In \emph{Formal Methods in Computer-Aided
Design (FMCAD)}, LNCS 1954, 127--144.
\end{thebibliography}

\vspace{1em}
\noindent\emph{Correspondence: Manuel G. Galmanus, Bluewave AI Research. This
is a preprint and has not been peer reviewed. All quantitative claims about the
compiler (verdicts, exit codes, bound, dependency count) are reproducible from
the open-source \texttt{axl-compiler} crate; the bound in Theorem~1 is
machine-discharged by an SMT solver and re-verifiable with \texttt{axlc
verify-cert}. The limitations in Section~\ref{sec:scope} are stated as
falsifiable conditions.}

\end{document}
